\section{Harness Ablations}\label{app:ablations}

This appendix gives the full per-component attribution behind the $\mathcal{H}_{\min}$-to-$\mathcal{H}_\mathrm{CH}$ progression gap (\cref{sec:exp:mechanisms}) and the reset-free bootstrap transfer results.

\subsection{Mechanism Attribution}\label{app:mechanisms}

\subsubsection{Pathfinding skills}\label{app:mechanisms:pathfinding}

\cref{fig:pathfinding} in the main text shows the Dijkstra-oracle ratio and cumulative skill calls. Below we give the measurement details and the residual structure the main-text summary omits.

For each first-traversal segment between consecutive milestones, we compute a BFS-optimal path length on the union of tiles observed by any run of that game on that map, and divide the agent's issued button presses within the segment by that optimum. Dialogue and battle presses are filtered out so the comparator only sees navigation. The ordering inverts inside gyms where dialogue and puzzle state dominate over navigation, which is why the residual gap to $\mathcal{H}_\mathrm{expert}$ in the main-text progression figure concentrates there. The refined library is heavily biased toward BFS and A$^*$ wrappers because saved presses on navigation translate directly into faster milestones, the strongest local signal the refinement loop sees. $\mathcal{H}_\mathrm{expert}$ routes navigation through a pre-built A$^*$ tool not visible to the \texttt{run\_skill} counter, which is why its curve sits at the floor in the main-text right panel. Most of the $\mathcal{H}_{\min}$-to-$\mathcal{H}_\mathrm{CH}$ delta is absorbed by the skill library alone.

\subsubsection{Skill debugging}\label{app:mechanisms:debug}

\begin{figure}[t]
\centering
\begin{minipage}[c]{0.68\linewidth}
  \centering
  \includegraphics[width=\linewidth]{h3_skill_timeline.pdf}
\end{minipage}\hfill
\begin{minipage}[c]{0.30\linewidth}
  \centering
  \includegraphics[width=\linewidth]{h3_skill_funnel.pdf}
\end{minipage}
\caption{Left: per-seed skill lifetime for the three Red from-scratch $\mathcal{H}_\mathrm{CH}$ runs. Markers are \texttt{add}, \texttt{update}, \texttt{run\_skill} (green if no re-update within 5 steps, red otherwise), and \texttt{delete}. Right: create-and-forget funnel across both games and both $\mathcal{H}_\mathrm{CH}$ bootstrap variants.}
\label{fig:skill_debug}
\end{figure}

For every persisted skill that sees at least one \texttt{skills\_updated} event in the refinement log, we compute the rolling success rate over a window of invocations immediately before and immediately after each update. We also track the create-to-succeed funnel: skills authored, invoked at all, invoked repeatedly, and ever successful. Reset-free refinement produces dramatic repairs on the skills the agent actually relies on, and the repair happens in the same episode where the failure occurred. \cref{fig:skill_debug} shows this with two complementary views. The left panel is per-seed skill lifetime: each skill occupies one lane, markers are add/update/run/delete events, and the update-to-next-execute edge is drawn so each debug iteration reads as an update followed by a run. The right panel is the create-and-forget funnel: most authored skills are never invoked, a small working set absorbs the bulk of calls, and even fewer see success. The refinement loop therefore triages: it repairs the skills the agent depends on, tolerates regressions on unused ones, and accepts a long create-and-forget tail. This is the argument for reset-free operation over reset-based baselines: the failure record and the repair sit inside the same trajectory, so the loop closes within a run rather than across resets.

\subsubsection{Sub-agent handoffs}\label{app:mechanisms:subagent}

\begin{figure}[t]
\centering
\includegraphics[width=\linewidth]{h4_subagent.pdf}
\caption{Sub-agent handoffs. (a) Cumulative approximate tokens by role (orchestrator solid, sub-agent dashed). (b) Cumulative \texttt{execute\_custom\_subagent} count. (c) Per-task-type handoff success: \textbf{exit} is the percent of spans ending via \texttt{return\_to\_orchestrator}; \textbf{focus} is the percent of returns where the orchestrator either pursued the pre-handoff objective or crossed a milestone within ten subsequent steps.}
\label{fig:subagent}
\end{figure}

We segment each run into spans of orchestrator execution and sub-agent execution, track approximate per-step input tokens from prompt length, and label each sub-agent span by task type. Sub-agent handoffs serve two roles for the harness: they keep per-step cost low by giving the sub-agent a narrow specialized context, and they let the orchestrator resume its prior objective after the sub-agent returns. \cref{fig:subagent} makes both points visible. The left panel plots orchestrator tokens and sub-agent tokens on a shared step axis; the sub-agent curve sits about an order of magnitude below the orchestrator curve throughout, which is the per-step saving the harness buys by partitioning context. The middle panel tracks cumulative handoff counts per condition; bootstrap-updating tracks from-scratch closely throughout the run, with a small plateau gap by run end. The right panel scores post-return behavior per task type; clean-return and on-task-recovery rates sit near the top of the scale for navigation, dialogue, and menu tasks. The harness rather than the raw model carries most of the long-horizon performance: once the orchestrator can delegate to cheap specialized contexts and trust the return, long tasks become tractable with far fewer tokens than the raw context would imply.

\subsubsection{Memory reuse}\label{app:mechanisms:memory}

\begin{figure}[t]
\centering
\includegraphics[width=\linewidth]{h8_memory_bottleneck.pdf}
\caption{\texttt{process\_memory} use inside the first real bottleneck of each game: Mauville Gym (Emerald, left column) and Mt Moon (Red, right column). One representative seed per condition. Top row: per-run timeline where each marker is a \texttt{process\_memory} tool call at the step it fired; color encodes provenance of the memory entry. Bottom row: the same ops aggregated into a stacked composition bar.}
\label{fig:memory}
\end{figure}

Every orchestrator step prompt lists the IDs and titles of all stored memories under a \texttt{LONG-TERM MEMORY OVERVIEW} section, so the agent sees the full catalog for free. The question is whether the agent \emph{pulls} on that catalog: requests the full content of an entry via \texttt{process\_memory}, invokes a memorized skill, or cites an entry ID in its reasoning. We measure this pull rate as the fraction of available entries ever referenced in an episode, and we localize reads to the milestone windows where the agent should need them: Mauville Gym on Emerald, Mt Moon on Red. Memory is leveraged once the library is both mature and inherited (\cref{fig:memory}): bootstrap runs, which load a from-scratch memory store at the start, consult it actively inside the gym and cave segments, while from-scratch runs write many entries and rarely reach back for them. The reference rate remains low in absolute terms, which we report honestly; most authored entries sit unused. The transferable unit of the framework is therefore the harness across runs, not a single episode, and an explicit reuse prior is a natural next step.

\subsection{Reset-Free Bootstrap Transfer}\label{app:bootstrap}

The bootstrap-updating variant tests whether the transferable unit is a single episode or the harness across episodes. On Emerald, every store the agent exercises during a bootstrap-updating run still targets the inherited harness. On Red, memory and skill invocations stay inherited, but the bootstrap-updating agent collapses its sub-agent budget to a handful of calls, and the few it makes cite IDs not present in the bootstrap. When that collapse happens, the milestone staircase regresses. The harness-as-transferable-unit claim therefore holds when the agent continues to exercise the inherited components and breaks when it abandons them: a reuse prior or a sub-agent deletion policy is the natural next step.

Bootstrap runs load the final skills, sub-agents, and memory of a successful from-scratch run before any new play. For every invocation we classify whether the invoked entry originated in the bootstrap or was authored during the bootstrap run itself, using retrieval semantics that count actual use across stores rather than passive prompt inclusion. Bootstrap succeeds by leveraging the inherited harness, and regressions show up where the agent stops using it. The end-of-run inherited share per store is reported in \cref{tab:c5_inheritance}. On Emerald every store that the agent exercises targets the inherited harness, with substantial per-run invocation counts on skills and sub-agents. On Red, memory and skill shares also stay inherited, and the only anomaly is sub-agents: the Red bootstrap-updating agents collapse their sub-agent budget to a handful of calls and the few calls they make cite IDs not present in the bootstrap. The milestone staircase in \cref{fig:progression} reads this phenotypically: Emerald bootstrap tracks from-scratch closely, while Red bootstrap-updating drifts below from-scratch and then below $\mathcal{H}_{\min}$ in lockstep with the collapse of sub-agent use. The harness rather than the episode is the transferable unit: reset-free refinement carries across runs when the inherited components stay in play, and it breaks when the agent abandons them.

% c5_inheritance.tex -- phase-2 inheritance rate (invocations originating in bootstrap)
\begin{table}[t]
\centering
\small
\caption{C5 inheritance: fraction of phase-2 invocations whose target was present in the phase-1 bootstrap. Mean across seeds (n=3 for each cell). ``--'' means no invocations of that store type in the run.}
\label{tab:c5_inheritance}
\begin{tabular}{llrr}
\toprule
Game & Store & Frozen & Continued \\
\midrule
Emerald & skills & 100.0\% $\pm$ 0.0 & 99.6\% $\pm$ 0.6 \\
 & subagents & 100.0\% $\pm$ 0.0 & 100.0\% $\pm$ 0.0 \\
 & memories & 98.2\% $\pm$ 1.9 & 100.0\% $\pm$ 0.0 \\
\midrule
Red & skills & 100.0\% $\pm$ 0.0 & 96.5\% $\pm$ 3.8 \\
 & subagents & 100.0\% $\pm$ 0.0 & 6.4\% $\pm$ 5.7 \\
 & memories & 100.0\% $\pm$ 0.0 & 100.0\% $\pm$ 0.0 \\
\bottomrule
\end{tabular}
\end{table}

\subsubsection{Red Bootstrap-Updating Regression}\label{app:red_regression}

The Red bootstrap-updating regression visible in \cref{fig:progression} begins around step 213. Newly authored sub-agents overtake the inherited ones at that point. They have not gone through the repair cycle observed for from-scratch skills (\cref{fig:skill_debug}), so their per-invocation success rate sits below what the bootstrap sub-agents reached before they were cached. A reuse prior on sub-agent selection, or a deletion rule that deprecates newly authored sub-agents whose task signature is covered by inherited ones, are natural follow-ups.
